\begin{sectionbox}
    \subsection{Fourierpolynom}
    \begin{description}
        \item[Periode] betragsmäßig kleinste Zahl, wo sich die Funktion $\forall t$ wiederholt,
              z.B. $f(t)=\cos(k\omega t) + i\sin(k\omega t)$ wiederholt sich für $|k|\omega T_k=2\pi \Rightarrow T_k=\frac{2\pi}{|k|\omega}$
        \item[T-periodisch] Falls $f(t + T)=f(t) \quad \forall t\in\mathbb R$
        \item[Kreisfrequenz] $\omega = \frac{2\pi}{T}$ mit Periode $T$
    \end{description}

    Sei $\omega > 0, c_k \in \mathbb C, -n \le k \le n, n \in \mathbb N$. Dann heißt
    \begin{equation*}
        S_n(t)=\sum_{k=-n}^{n}c_ke^{ik\omega t}, t\in\mathbb R
    \end{equation*}
    Fourierpolynom bzw. trigonometrisches Polynom.

    \subsubsection{Eigenschaften}
    \begin{itemize}
        \item $T_1$ ist die größte auftretende Periode aller Summanden $\Rightarrow S_n$ ist $T_1$-periodisch
        \item $T_k= \frac{T_1}{|k|}$
        \item $\omega$ ist für alle Summanden gleich, $k$ variiert
    \end{itemize}
\end{sectionbox}

\subsection{Fourierreihe}
Jedes Fourierpolynom ist gleichzeitig eine Fourierreihe, wenn man alle nicht benötigten $c_k=0$ setzt.
\begin{emphbox}
    \begin{equation*}
        S(t)=\lim_{n\rightarrow\infty}S_n(t)=\sum_{k\in\mathbb Z}c_ke^{ik\omega t}=\sum_{k=-\infty}^{\infty}c_ke^{ik\omega t}
    \end{equation*}
\end{emphbox}

\begin{sectionbox}
    \subsubsection{Bestimmung der Fourierkoeffizienten}
    Sei eine $T$-periodische Funktion $f: \mathbb R \rightarrow \mathbb C$ gegeben und $k\in\mathbb Z$. Dann heißt
    \begin{equation*}
        c_k=\frac{1}{T}\int_a^{a+T} f(t)e^{-ik\omega t}\diff t, \quad a \text{ bel.}
    \end{equation*}
    der $k$-te Fourierkoeffizient von $f$. \\
    \textbf{Fourierspektrum} Familie $\{c_k\}_{k\in\mathbb Z}$ aller Fourierkoeffizienten. Beschreibt Amplituden und Phasenverschiebungen der Teilschwingungen in $f$.

    \subsubsection{Konvergenz der Fourierreihe}
    \begin{sectionbox}
        $f:\mathbb R \rightarrow \mathbb C$ ist stückweise stetig, falls $\forall t \in \mathbb R$ die rechts- und linksseitigen Grenzwerte existieren und $f$ auf jedem beschränkten Interval nur in endlich vielen Punkten unstetig ist (Sprungstellen)
    \end{sectionbox}

    \begin{sectionbox}
        $f$ stückweise stetig ist stückweise stetig differenzierbar, falls $f$ auf jedem beschränkten Intervall nur in endlich vielen Punkten nicht stetig differenzierbar ist und die rechts- und linksseitigen Grenzwerte von $f'$ auch existieren
    \end{sectionbox}

    \begin{description}
        \item[Punktweise Konvergenz] Fourierreihe konvergiert $\forall t$ gegen $\frac{f(t^-) + f(t^+)}{2}$
        \item[Gleichmäßige Konvergenz] Auf jedem abgeschlossenen Intervall ohne Sprungstellen konvergiert die Fourierreihe gegen $f$
    \end{description}
\end{sectionbox}

\begin{sectionbox}
    \subsubsection{Orthogonalreihenrelationen}
    Für $k_1, k_2 \in \mathbb Z$ gilt:
    \begin{equation*}
        \frac{1}{T}\int_0^Te^{ik_1\omega t}e^{-ik_2\omega t}\diff t = \delta_{k_1,k_2}=\begin{cases}
            1, & k_1=k_2    \\
            0, & k_1\ne k_2
        \end{cases}
    \end{equation*}
    $\Rightarrow f_k(t)=e^{ik\omega t}$ hat Fourierspektrum $c_{k'}=\delta_{k,k'} \; \forall k' \in \mathbb Z$
\end{sectionbox}